\section{Kesimpulan}
Milestone 4 berhasil menambahkan tahap intermediate code generation dan interpreter pada Arion Compiler. Program dapat menerima source code, token stream, atau parse tree, kemudian menghasilkan parse tree, AST, symbol table, decorated AST, intermediate code, dan output eksekusi program. Implementasi intermediate code generator telah mengubah decorated AST menjadi daftar instruksi sekuensial yang terdiri dari \texttt{LIT}, \texttt{LOD}, \texttt{STO}, \texttt{INT}, \texttt{JMP}, \texttt{JPC}, \texttt{OPR}, \texttt{CAL}, dan \texttt{RET}. Instruksi-instruksi tersebut kemudian dieksekusi oleh interpreter yang berperan sebagai virtual machine dengan siklus fetch-decode-execute.

Stack machine telah diimplementasikan untuk mengelola memori eksekusi melalui stack frame yang menyimpan variabel lokal, argumen, static link, dynamic link, dan return address. Interpreter mampu menangani translasi control flow dari \texttt{if}, \texttt{while}, \texttt{for}, \texttt{repeat}, dan \texttt{case} ke dalam instruksi sekuensial dengan label dan jump. Program juga telah mendukung array access dengan satu dan dua dimensi melalui instruksi \texttt{ALOD} dan \texttt{ASTO}.

Berdasarkan pengujian, program berhasil menangani kasus dasar, control flow bersarang, perulangan, array, dan runtime error. Program juga mampu melaporkan beberapa runtime error secara informatif tanpa crash, termasuk division by zero, stack overflow, array out of bounds, invalid jump targets, dan integer overflow.

\section{Saran}
Pengembangan berikutnya dapat menambahkan dukungan untuk function call dalam ekspresi (\texttt{FuncCallNode}) pada level intermediate code generation, sehingga nilai return dari fungsi dapat digunakan dalam assignment atau ekspresi lainnya. Selain itu, record field access dapat diperluas agar field dapat diakses dan dimodifikasi pada saat runtime. Boolean operator \texttt{and}, \texttt{or}, dan \texttt{not} juga dapat ditambahkan ke instruction set agar kondisi kompleks dapat dieksekusi.

Untuk meningkatkan keamanan, interpreter dapat ditambahkan dengan deteksi use-after-free pada variabel yang sudah di-pop dari stack, serta validasi tipe runtime untuk memastikan operasi hanya dilakukan pada tipe data yang kompatibel. Pengujian juga dapat diperluas dengan kasus edge case yang lebih banyak, termasuk nested procedure, array multidimensi lebih dari dua, dan kombinasi control flow yang lebih kompleks.
